\section{Diagnostic Analysis}

\subsection{Retrieval Is the Bottleneck}
% 中文注释：这里改成更短的表达，同时顺手消除开头那条轻微的断行 warning。
We sampled 20 FAISS failures and grouped them by retrieval-context length.
% 中文注释：这里只统一符号间距和普通字体，不改变分析结论。
All 20 had short or empty retrieval context ($< 200$ characters), compared with an average of $\sim 500$ characters for correct answers.
This heuristic analysis attributes roughly $\sim 55\%$ of the sampled failures to retrieval and $\sim 45\%$ to downstream reasoning, suggesting that improvements in retrieval quality, particularly temporal-aware retrieval, may yield larger gains than improvements in reasoning alone.

\subsection{Confidence Is Anti-Calibrated}
% 中文注释：保留术语，但取消不必要的斜体强调，让正文更接近 ACL 常见风格。
A notable diagnostic finding is that some vector-based systems exhibit anti-calibrated confidence scores.
Table~\ref{tab:calibration} shows Expected Calibration Error (ECE), Brier score, and AUROC for each system.

\begin{table}[t]
\centering
\small
\begin{tabular}{lccc}
\toprule
\textbf{System} & \textbf{ECE}$\downarrow$ & \textbf{Brier}$\downarrow$ & \textbf{AUROC} \\
\midrule
FAISS & 0.163 & 0.229 & 0.453 \\
Naive RAG & 0.319 & 0.303 & 0.552 \\
BM25 & 0.235 & 0.769 & 0.467 \\
Simple KG & 0.672 & 0.672 & 0.500 \\
ChromaDB & 0.412 & 0.412 & 0.500 \\
\bottomrule
\end{tabular}
\caption{Confidence calibration metrics. FAISS and BM25 show AUROC $< 0.5$, meaning that wrong answers receive higher confidence than correct ones.}
\label{tab:calibration}
\end{table}

FAISS's AUROC of 0.453 indicates that confidence scores are worse than random at predicting correctness: wrong answers often receive higher confidence than correct ones.
This has practical implications, because agents cannot reliably use internal confidence alone to determine whether a temporal-memory answer should be trusted or externally verified.

\subsection{Knowledge Retraction vs.\ Version Tracking}
We decompose performance on questions involving knowledge dynamics (Table~\ref{tab:retraction}).
Simple KG shows an asymmetric profile: it is relatively strong on version tracking (81.0\%, close to FAISS's 82.8\%) but weaker on knowledge retraction (53.6\%).
This pattern suggests that graph structures may naturally support tracking fact evolution, while handling explicit revocation of superseded knowledge remains more difficult.

\begin{table}[t]
\centering
% 中文注释：该表继续保持按栏宽缩放，只统一字体和符号间距，不改变尺寸策略。
\resizebox{\columnwidth}{!}{%
\begin{tabular}{lccc}
\toprule
\textbf{System} & \textbf{Retraction (28Q)} & \textbf{Ver.\ Track.\ (58Q)} & \textbf{Neither (222Q)} \\
\midrule
FAISS & 67.9\% & \textbf{82.8\%} & 74.8\% \\
Naive RAG & 67.9\% & 75.9\% & 74.3\% \\
BM25 & 57.1\% & 74.1\% & \textbf{75.2\%} \\
Simple KG & 53.6\% & 81.0\% & 65.3\% \\
ChromaDB & 39.3\% & 37.9\% & 42.3\% \\
\bottomrule
\end{tabular}%
}
\caption{Accuracy by knowledge-dynamics type. Simple KG is strong at version tracking but weak at retraction. Bootstrap CIs for retraction span 35pp because $n = 28$ is small.}
\label{tab:retraction}
\end{table}

\subsection{Implications for Memory System Design}
Our findings suggest three design implications:

\paragraph{1. Temporal-aware retrieval.}
Since $\sim 55\%$ of sampled failures are retrieval-limited, adding temporal metadata filtering (e.g., ``retrieve only facts recorded before $T$'') or bitemporal indexing may yield larger gains than improving the reasoning model alone.

\paragraph{2. Hybrid vector+graph architectures.}
% 中文注释：这里只统一字体，不改原始推断力度。
The complementarity between systems (oracle +9.8pp) and Simple KG's adversarial robustness suggest that combining dense retrieval's breadth with graph structure's precision could outperform either approach alone.

\paragraph{3. External verification for temporal claims.}
Anti-calibrated confidence means that internal signals are unreliable for identifying uncertain temporal answers. External verification, such as cross-checking against source timestamps, may therefore be especially valuable for bitemporal queries.
